%%%%%%%%%%  scan idx 192  |  folio 179  %%%%%%%%%%
\textbf{Remarque 2.4.1.1.} Si $S$ a «~assez de points fermés~»\footnote{i.e. assez de foncteurs fibres (SGA~4 IV 6.4.1) du type $i_s:s\to S$, où $i_s$ est l'inclusion d'un fermé de $S$.} $s$ tels que l'Idéal $\mathfrak m_s$ défini par $s$\footnote{i.e. formé des sections de $\mathcal O_S$ dont la fibre en $s$ appartient à l'idéal maximal $m_s$ de $\mathcal O_{S,s}$.} soit cohérent, alors la relation «~$H^1(S,F)=0$ pour tout faisceau cohérent $F$~» implique le théorème A. En effet, soient $F$ un $\mathcal O_S$-Module cohérent et $s$ un point fermé de $S$. On a la suite exacte
\[
0\longrightarrow\mathfrak m_sF\longrightarrow F\longrightarrow F/\mathfrak m_sF\longrightarrow0,
\]
où $F/\mathfrak m_sF$ est un faisceau cohérent concentré sur $s$. Appliquant le foncteur $H^0(S,.)$ et écrivant la suite exacte de cohomologie, on trouve que $H^0(S,F)\longrightarrow H^0(S,F/\mathfrak m_sF)$ est un épimorphisme. Comme on a $H^0(S,F/\mathfrak m_sF)\xrightarrow{\sim}F_s/m_sF_s$, où $m_s$ est l'idéal maximal de $\mathcal O_{S,s}$, on conclut par Nakayama.

Le rédacteur ignore si l'implication précédente est vraie en général (quand on suppose seulement que $\mathcal O_S$ est un faisceau cohérent d'anneaux locaux).

Le rédacteur ignore également si la relation «~$H^1(S,F)=0$ pour tout faisceau cohérent $F$~» implique le théorème B.

\textbf{Exemples 2.4.2.} Un espace analytique complexe de Stein [3], un compact convexe de $\mathbb C^n$ muni du faisceau induit par le faisceau des fonctions holomorphes [3], le spectre d'un anneau noethérien, un espace rigide analytique affine [13] définissent des topos de Stein.

\textbf{Lemme 2.4.3.} Soit $S$ un topos de Stein (2.4.1). On suppose l'objet final $S$ quasi-compact (2.1.1), et l'on note $\Coh(S)$ la catégorie abélienne des $\mathcal O_S$-Modules cohérents, i.e. de présentation finie. Alors :
\begin{enumerate}[label=\alph*)]
\item La catégorie $\Lib(S)$ (resp. $\Loclib(S)$) (2.0) est relevable dans $\Coh(S)$ et quasi-stable (resp. stable) par noyau d'épimorphisme (I 1.2).
\item Les foncteurs canoniques
\[
\mathrm K^{-}(\Lib(S))\longrightarrow\mathrm K^{-}(\Loclib(S))\longrightarrow\mathrm D^{-}(\Coh(S))
\]
sont des équivalences de catégories.
